\documentclass[11pt]{article}
\usepackage[a4paper,margin=25mm]{geometry}
\usepackage{amsmath,amssymb,amsthm,booktabs,array}
\newtheorem{theorem}{Theorem}\newtheorem{lemma}{Lemma}
\title{Literal structural descriptions in fourteen retained $L_2(q)$ tables}
\author{Private recovery component}\date{}
\begin{document}\maketitle
\begin{theorem}
For the exact current companion tables and the accepted numerical row
bijections listed below, all 348 displayed structural descriptions hold for
their actual retained subgroup representatives. This asserts the indicated
decompositions, not that an extension label determines a unique group.
\end{theorem}
\begin{center}\begin{tabular}{rrrrrrr}\toprule
$q$&7&8&11&13&17&19\\
rows&15&12&16&16&22&19\\\midrule
$q$&16&23&25&27&29&31\\
rows&21&23&37&16&22&29\\\midrule
$q$&32&64&&&&\\
rows&24&76&&&&\\\bottomrule
\end{tabular}\end{center}
The tables are attributed to Thomas Connor and Dimitri Leemans. All nineteen
native pages were inspected. Their dates range from November 2012 to June
2013, but identity with every historical file available to the 2013 source
article is not claimed. The retained numerical classifications and weighted
incidence checks are separate acquired components. The term ``lattice'' in
the source headings does not supply a lattice property here.

The token trees in \texttt{native\_assertions.json} preserve the literal
punctuation and left association: a colon denotes the displayed split
extension, a centered dot the displayed nonsplit extension, and $p^e$ an
elementary abelian group. Bare integers denote cyclic groups. In particular,
row 5 of $L_2(32)$ is cyclic of order 33, whereas row 4 of $L_2(27)$ is
$3^3$. The dihedral subscript denotes group order. No exponent is recovered
by guessing how to split a flattened decimal string.

\paragraph{Retained inputs.}
The acquired complete inventories supply actual permutation generators for
all subgroup representatives and the exact weighted row bijections. The
robust cached-weight calculation supplies the final orders and incidences;
no previous local-class multiplicity assumption is reintroduced. The fourteen
top rows use the original executed $\operatorname{PSL}(2,q)$ constructors,
their literal saved ambient models and the full-order subgroup in each
accepted inventory. A subgroup of a finite group having the full group order
equals that group. Its selected generators need not be the constructor's
generators. This route retains the named-constructor identification premise;
order alone is never used to identify an arbitrary group with $L_2(q)$.

The registry acquires the 29 exact small models and decomposition certificates
from the accepted A5/M11 structural components. It also uses the literal
retained $L_2(8)$ model for the proper occurrence in $L_2(64)$. The latter
occurrence is paid by an explicit isomorphism, not by an ambient name.
The generator-edge transport theorem MC.3 in the accepted marked-coordinate
umbrella is acquired in full. No ambient subgroup census is run here.

\newpage
\paragraph{Efficient literal transport.}
All permutations act on the right: $ab$ sends $i$ first to $i^a$ and then to
$(i^a)^b$. Each concrete registry model has ordered generators and a complete
positive-word closure. A stored vector $f$ is checked to map identity to
identity, to be a bijection onto the closure of the actual target generators,
and to satisfy $f(xa)=f(x)f(a)$ on every right-generator edge. Induction on
words gives the unique homomorphism; bijectivity makes it an isomorphism.
This is precisely acquired MC.3. The cost is proportional to the number of
group elements times the number of generators, not to a full multiplication
table. New models are indexed by their literal data, never by a supposedly
canonical extension name. Factor references point only to previously proved
models, so the certificate dependency order is acyclic.

\paragraph{Atoms and split extensions.}
A cyclic atom has an element of the asserted order in a literal group of
that order. An elementary abelian atom has the asserted number of commuting
order-$p$ generators and their complete closure has order $p^e$. A dihedral
atom has $r,s$ with orders $m,2$, $srs=r^{-1}$, and their closure has order
$2m$. These conditions prove the displayed atom. The $A_4,A_5,S_3$ atoms and
the already certified small split types use the acquired registry.

At a new split node $H=N:K$, the verifier constructs $N$ and $K$ from their
actual generators, checks membership in $H$, and transports their recursively
certified descriptions. Conjugation of each generator of $N$ by each generator
of $H$ remains in $N$. Finiteness makes these inclusions equalities and proves
$N\trianglelefteq H$. It checks $N\cap K=1$ and $|N||K|=|H|$; hence $H=NK$
and $H\cong N\rtimes K$. A noncommuting generator pair additionally pays the
displayed nontrivial action. This does not assert that all actions with the
same factors are isomorphic. Every nested node follows the exact source tree.

\begin{lemma}
Let $N\trianglelefteq H$ have prime index $p$. The quotient map onto $H/N$
splits if and only if some element outside $N$ has order $p$.
\end{lemma}
\begin{proof}
A complement has order $p$, so any of its nonidentity elements gives such an
element. Conversely, an outside element of order $p$ generates a subgroup
meeting $N$ trivially and mapping onto the quotient. It is a complement.
\end{proof}
For each centered-dot node in this cluster, the quotient has order 2 or 3.
The actual normal subgroup and its complete factor certificate are checked
as above. The verifier then checks $x^p\ne1$ for \emph{every} element of
$H\setminus N$. This proves nonsplitting by the lemma. Failure to find a
complement is not used as evidence. The prime-three case is the elementary
extension of the previously accepted prime-two criterion.

\newpage
\paragraph{Complete certificate and controls.}
The deterministic literal verifier pays all 334 proper-subgroup occurrences,
with 70 new concrete registry models in addition to the 30 acquired model
inputs. There are 387 occurrence/factor transports, 69,818 right-generator
edge checks and 3,488 outside-element checks for nonsplitting. Together with
the fourteen exact top-model routes, these account for all 348 rows. The
source join checks every original generator array and numerical row map;
none of the source names filtered an ambient group search.

Candidate generation only sought witnesses inside the retained proper
subgroups. It completed in 3.7 seconds under a fixed total 120-second budget.
GAP normal-subgroup and isomorphism routines suggested these witnesses;
their search completeness is not a premise of the literal certificate. A
startup failure caused by an empty generator list was preserved before the
successful pass. One successful GAP output split an integer by its standard
backslash-newline continuation. The raw file is preserved and the decoder
permits removal only when both adjacent characters are decimal digits.
No mathematical search was repeated for that formatting issue.

The controls damage identity images, a bijective nonhomomorphism, element
coverage, source-row coverage, actual generator and row-map joins, source
superscripts and punctuation, quotient orders, elementary bases, dihedral
relations and the factor dependency order. All must be rejected. Separate
literal controls distinguish the nonsplit cyclic groups of orders 4 and 9
from split elementary abelian groups of those orders. Ordinary and optimized
Python execute the same explicit checks; assertions are not used as gates.

\paragraph{Scope and attribution.}
This is a finite structural-label recovery on the fourteen specified current
tables, conditional on the already acquired complete inventories and explicit
named-constructor inputs. It proves neither a new classification of subgroups
of $L_2(q)$ nor the identities of unexamined historical editions. It does not
complete the entire subgroup atlas or add a new original-source, conjecture,
formalization or measurement credit. The normal-series, semidirect-product
and generator-edge arguments are standard, with their existing project
interfaces retained. Root suggested the cluster/registry strategy and reuse
of MC.3; the author prepared the native ledger, witnesses and literal checker.
\end{document}
